% Figure from MATH 590 notes, section 11. Compile with the notes preamble.
\begin{tikzpicture}[scale=1.5]
  \draw[figR, thick] (-1,-1) rectangle (1,1);
  \fill[figB!12] (0,0) circle (1);
  \draw[figB, thick] (0,0) circle (1);
  \fill[figR!15] (-0.7071,-0.7071) rectangle (0.7071,0.7071);
  \draw[figR, thick] (-0.7071,-0.7071) rectangle (0.7071,0.7071);
  \fill[black] (0,0) circle (1.1pt) node[below left=0pt, font=\scriptsize] {$x$};
  \draw[black!55, thin, -stealth] (0,0) -- (0,-1) node[pos=0.72, right=-1pt, font=\scriptsize, black!70] {$\varepsilon$};
  \draw[black!55, thin, -stealth] (0,0) -- (0.7071,0) node[pos=0.42, above=0pt, font=\scriptsize, black!70] {$\varepsilon/\sqrt2$};
  \node[font=\scriptsize, figR, anchor=west] at (1.05,0.95) {$B_\rho(x,\varepsilon)$};
  \node[font=\scriptsize, figB, anchor=west] at (1.05,0.35) {$B_d(x,\varepsilon)$};
  \draw[figB, thin] (1.05,0.35) -- (0.94,0.34);
  \node[font=\scriptsize, figR, anchor=west] at (1.05,-0.35) {$B_\rho(x,\varepsilon/\sqrt2)$};
  \draw[figR, thin] (1.05,-0.35) -- (0.6,-0.45);
\end{tikzpicture}
